Este apêndice complementa a Seção~\ref{cap:resultados:sec-analise-base-dados} com o detalhamento numérico do conjunto histórico e das 40 instâncias (10 por horizonte) utilizadas nos experimentos. A organização espelha a do capítulo: a Seção~\ref{apendice:base-geral-detalhe} aprofunda a análise descritiva do conjunto completo (02/01/2022 a 31/07/2023), e as Seções~\ref{apendice:base-h15}--\ref{apendice:base-h90} apresentam, para cada horizonte de planejamento, a tabela estendida das 10 instâncias, com colunas adicionais (médias diárias, número de unidades ativas) que foram omitidas do capítulo por concisão.

\noindent\textbf{Estimativa.}\, Como na Seção~\ref{cap:resultados:base-geral}, todas as razões \emph{Demanda / (Salas\,$\times$\,Turnos)} e \emph{Disponibilidade / Demanda} reportadas neste apêndice são \emph{limitantes ingênuos}. Elas \textbf{ignoram} as restrições \ref{eq:r1}--\ref{eq:r5} do modelo (qualificação médico--unidade, unicidade diária, capacidade por turno, cobertura mínima) e \textbf{não correspondem a uma alocação factível}; servem apenas como referência rápida para o regime oferta\,$\times$\,demanda esperado em cada recorte.

\section{Conjunto Histórico Completo}
\label{apendice:base-geral-detalhe}

A Figura~\ref{fig:apendice-scatter-salas-outros} examina a relação entre a soma mensal de salas disponíveis e três indicadores agregados na mesma granularidade: a demanda mensal de exames, a disponibilidade mensal de médicos e a quantidade de médicos distintos com pelo menos um turno ativo no mês. Demanda e disponibilidade exibem associação aproximadamente linear com a oferta de salas, enquanto a contagem de médicos distintos cresce de forma menos regular: o número de profissionais cadastrados aumenta ao longo do período sem que isso se traduza em proporção direta de salas adicionais, indicando que parte do crescimento da força de trabalho é absorvida pelo aumento da utilização das salas existentes.

\newpage

\begin{figure}[h]
	\centering
	\begin{subfigure}{0.333\textwidth}
		\includegraphics[width=\textwidth]{images/scatter-sala-demanda_2.png}
		\caption{Salas $\times$ Demanda}
		\label{fig:apendice-scatter-sala-dem}
	\end{subfigure}
	\begin{subfigure}{0.333\textwidth}
		\includegraphics[width=\textwidth]{images/scatter-sala-disponibilidade.png}
		\caption{Salas $\times$ Disponibilidade}
		\label{fig:apendice-scatter-sala-disp}
	\end{subfigure}
	\begin{subfigure}{0.333\textwidth}
		\includegraphics[width=\textwidth]{images/scatter-sala-coutmed.png}
		\caption{Salas $\times$ Médicos distintos}
		\label{fig:apendice-scatter-sala-med}
	\end{subfigure}
	\caption{Figuras de dispersão entre indicadores agregados por mês.}
	\label{fig:apendice-scatter-salas-outros}
	\fautor
\end{figure}

A distribuição de salas entre as 37 unidades clínicas, ilustrada na Figura~\ref{fig:apendice-bar-sala-unidade}, é fortemente assimétrica: poucas unidades concentram a maior parte da capacidade agregada de salas-turno, enquanto a cauda longa contém clínicas com participação marginal. Essa heterogeneidade é uma das fontes de dificuldade do modelo MIP: clínicas de grande porte impõem restrições de cobertura mínima (C5) sobre um volume substancial de exames, e a mobilidade dos médicos entre essas unidades (qualificação dada pela matriz $I \times U$) é frequentemente limitada.

\newpage

\begin{figure}[h]
	\centering
	\includegraphics[width=0.8\textwidth]{images/bar-sala-unidade.png}
	\caption{Soma de salas $\times$ turnos por unidade clínica no conjunto histórico (2022--2023).}
	\label{fig:apendice-bar-sala-unidade}
	\fautor
\end{figure}

A Figura~\ref{fig:apendice-histogramas} apresenta as distribuições marginais da demanda e da quantidade de salas disponíveis ao nível mais granular do conjunto (uma observação por tripla data\,$\times$\,turno\,$\times$\,unidade). A demanda exibe distribuição fortemente assimétrica à direita, com massa concentrada em valores baixos a moderados e cauda longa para unidades de grande porte; a quantidade de salas disponíveis concentra-se em 1, 2 e 3 salas por turno, refletindo a estrutura modular das unidades.

\begin{figure}[h]
	\centering
	\begin{subfigure}{0.45\textwidth}
		\includegraphics[width=\textwidth]{images/hist-demanda.png}
		\caption{Demanda por (data, turno, unidade)}
		\label{fig:apendice-hist-demanda}
	\end{subfigure}
	\begin{subfigure}{0.45\textwidth}
		\includegraphics[width=\textwidth]{images/hist-salas.png}
		\caption{Salas disponíveis por (data, turno, unidade)}
		\label{fig:apendice-hist-salas}
	\end{subfigure}
	\caption{Distribuições marginais de demanda e capacidade de salas no conjunto histórico (2022--2023).}
	\label{fig:apendice-histogramas}
	\fautor
\end{figure}

\clearpage

\section{Horizonte de 15 dias}
\label{apendice:base-h15}

O Quadro~\ref{quad:apendice-base-h15}, ao final desta seção, apresenta a caracterização estendida das 10 instâncias quinzenais, ordenadas cronologicamente. Em relação ao Quadro~\ref{quad:base-h15} do capítulo, são acrescentadas as colunas de médias diárias (demanda, disponibilidade e salas\,$\times$\,turnos) e o número de unidades clínicas ativas no recorte.

\noindent No horizonte quinzenal, a soma de exames demandados oscila entre $\sim$28k (recorte de janeiro/2022) e $\sim$44k (recortes de 2023), refletindo o crescimento gradual do volume operacional ao longo do período. A razão \emph{Demanda / (Salas\,$\times$\,Turnos)} permanece em uma faixa estreita (20{,}8--23{,}5), sem ruptura clara entre instâncias de parametrização e de validação. A razão \emph{Disponibilidade / Demanda} inicia em $\sim$1{,}08 nos primeiros recortes de 2022 (excesso de mão-de-obra disponível) e cai para $\sim$0{,}98 nos recortes de meados de 2023 (déficit incipiente), evidenciando que o crescimento da demanda nos últimos meses do conjunto histórico supera o crescimento da disponibilidade médica registrada na base.

\begingroup\let\cleardoublepage\clearpage
\begin{sidewaysquadro}[p]
	\centering
	\caption{Caracterização estendida das 10 instâncias do horizonte de 15 dias.}
	\begin{tabular}{|M{2.6cm}|c|c|c|c|c|c|c|c|c|c|}
\hline
\textbf{Inst\^ancia} & \textbf{M\'ed.} & \textbf{Unid.} & \textbf{$\Sigma$ Dem.} & \textbf{Dem./dia} & \textbf{$\Sigma$ Disp.} & \textbf{Disp./dia} & \textbf{$\Sigma$ Salas$\times$T} & \textbf{Salas/dia} & \textbf{Dem./Oferta} & \textbf{Disp./Dem.} \\
\hline\hline
02/01/22 a 16/01/22 & 127 & 35 & 27.855 & 1857 & 29.950 & 1997 & 1.338 & 89{,}2 & 20{,}82 & 1{,}08 \\
\hline
01/03/22 a 15/03/22 & 136 & 36 & 36.397 & 2426 & 37.374 & 2492 & 1.575 & 105{,}0 & 23{,}11 & 1{,}03 \\
\hline
01/05/22 a 15/05/22 & 145 & 36 & 36.358 & 2424 & 36.726 & 2448 & 1.548 & 103{,}2 & 23{,}49 & 1{,}01 \\
\hline
01/07/22 a 15/07/22 & 143 & 36 & 35.136 & 2342 & 37.459 & 2497 & 1.648 & 109{,}9 & 21{,}32 & 1{,}07 \\
\hline
01/09/22 a 15/09/22 & 139 & 37 & 36.467 & 2431 & 38.108 & 2541 & 1.691 & 112{,}7 & 21{,}57 & 1{,}04 \\
\hline
01/11/22 a 15/11/22 & 147 & 37 & 36.316 & 2421 & 37.454 & 2497 & 1.609 & 107{,}3 & 22{,}57 & 1{,}03 \\
\hline
01/01/23 a 15/01/23 & 150 & 37 & 31.350 & 2090 & 32.371 & 2158 & 1.416 & 94{,}4 & 22{,}14 & 1{,}03 \\
\hline
01/03/23 a 15/03/23 & 162 & 37 & 44.210 & 2947 & 43.319 & 2888 & 1.920 & 128{,}0 & 23{,}03 & 0{,}98 \\
\hline
01/05/23 a 15/05/23 & 171 & 37 & 43.904 & 2927 & 43.193 & 2880 & 1.895 & 126{,}3 & 23{,}17 & 0{,}98 \\
\hline
01/07/23 a 15/07/23 & 171 & 37 & 44.332 & 2955 & 44.019 & 2935 & 2.046 & 136{,}4 & 21{,}67 & 0{,}99 \\
\hline
\end{tabular}

	\label{quad:apendice-base-h15}
	\fautor
\end{sidewaysquadro}
\endgroup

\clearpage

\section{Horizonte de 30 dias}
\label{apendice:base-h30}

O Quadro~\ref{quad:apendice-base-h30}, ao final desta seção, apresenta a caracterização estendida das 10 instâncias mensais. Em relação ao Quadro~\ref{quad:base-h30} do capítulo, são acrescentadas as colunas de médias diárias (demanda, disponibilidade e salas\,$\times$\,turnos) e o número de unidades clínicas ativas no recorte.

\noindent No horizonte mensal, a soma de exames demandados varia entre $\sim$57k (janeiro/2022, instância parcial por iniciar em 02/01) e $\sim$91k (março/2023). A média diária de demanda evolui de $\sim$1.918 para $\sim$3.031 exames, com aumento simultâneo no número de médicos distintos ativos (141 para 182). A razão \emph{Demanda / (Salas\,$\times$\,Turnos)} oscila entre 21{,}5 e 23{,}6, e a razão \emph{Disponibilidade / Demanda} segue o mesmo padrão observado no horizonte quinzenal: começa próxima de 1{,}07 em 2022 e converge para valores ligeiramente inferiores a 1 nos recortes de 2023. A janela mensal já é suficiente para suavizar parte da variabilidade semanal observada no horizonte de 15 dias, mas preserva integralmente o efeito sazonal anual.

\begingroup\let\cleardoublepage\clearpage
\begin{sidewaysquadro}[p]
	\centering
	\caption{Caracterização estendida das 10 instâncias do horizonte de 30 dias.}
	\begin{tabular}{|M{2.6cm}|c|c|c|c|c|c|c|c|c|c|}
\hline
\textbf{Inst\^ancia} & \textbf{M\'ed.} & \textbf{Unid.} & \textbf{$\Sigma$ Dem.} & \textbf{Dem./dia} & \textbf{$\Sigma$ Disp.} & \textbf{Disp./dia} & \textbf{$\Sigma$ Salas$\times$T} & \textbf{Salas/dia} & \textbf{Dem./Oferta} & \textbf{Disp./Dem.} \\
\hline\hline
02/01/22 a 31/01/22 & 141 & 35 & 57.553 & 1918 & 61.723 & 2057 & 2.676 & 89{,}2 & 21{,}51 & 1{,}07 \\
\hline
01/03/22 a 30/03/22 & 144 & 36 & 74.772 & 2492 & 76.578 & 2553 & 3.197 & 106{,}6 & 23{,}39 & 1{,}02 \\
\hline
01/05/22 a 30/05/22 & 154 & 36 & 74.067 & 2469 & 75.659 & 2522 & 3.218 & 107{,}3 & 23{,}02 & 1{,}02 \\
\hline
01/07/22 a 30/07/22 & 157 & 36 & 70.942 & 2365 & 74.545 & 2485 & 3.279 & 109{,}3 & 21{,}64 & 1{,}05 \\
\hline
01/09/22 a 30/09/22 & 146 & 37 & 73.809 & 2460 & 76.854 & 2562 & 3.367 & 112{,}2 & 21{,}92 & 1{,}04 \\
\hline
01/11/22 a 30/11/22 & 154 & 37 & 75.006 & 2500 & 76.229 & 2541 & 3.266 & 108{,}9 & 22{,}97 & 1{,}02 \\
\hline
01/01/23 a 30/01/23 & 163 & 37 & 70.932 & 2364 & 72.476 & 2416 & 3.146 & 104{,}9 & 22{,}55 & 1{,}02 \\
\hline
01/03/23 a 30/03/23 & 175 & 37 & 90.943 & 3031 & 88.115 & 2937 & 3.862 & 128{,}7 & 23{,}55 & 0{,}97 \\
\hline
01/05/23 a 30/05/23 & 178 & 37 & 90.052 & 3002 & 87.432 & 2914 & 3.830 & 127{,}7 & 23{,}51 & 0{,}97 \\
\hline
01/07/23 a 30/07/23 & 182 & 37 & 88.366 & 2946 & 87.615 & 2920 & 3.979 & 132{,}6 & 22{,}21 & 0{,}99 \\
\hline
\end{tabular}

	\label{quad:apendice-base-h30}
	\fautor
\end{sidewaysquadro}
\endgroup

\clearpage

\section{Horizonte de 60 dias}
\label{apendice:base-h60}

O Quadro~\ref{quad:apendice-base-h60}, ao final desta seção, apresenta a caracterização estendida das 10 instâncias bimestrais. Em relação ao Quadro~\ref{quad:base-h60} do capítulo, são acrescentadas as colunas de médias diárias (demanda, disponibilidade e salas\,$\times$\,turnos) e o número de unidades clínicas ativas no recorte.

\noindent No horizonte bimestral, a soma de exames demandados ultrapassa a casa das centenas de milhares (de $\sim$130k a $\sim$180k por instância), com média diária estável na faixa $\sim$2.160--2.993 exames. O número de médicos distintos cresce de 146 (recorte iniciado em 02/01/2022) para 193 (recortes de meados de 2023), confirmando a tendência de expansão da força de trabalho registrada na base. A amplitude das duas razões \emph{oferta\,$\times$\,demanda} se reduz em relação aos horizontes mais curtos --- \emph{Demanda / (Salas\,$\times$\,Turnos)} concentra-se entre 21{,}6 e 23{,}7, e \emph{Disponibilidade / Demanda} entre 0{,}96 e 1{,}06 ---, indicando que janelas mais longas amortecem flutuações pontuais de demanda e disponibilidade.

\begingroup\let\cleardoublepage\clearpage
\begin{sidewaysquadro}[p]
	\centering
	\caption{Caracterização estendida das 10 instâncias do horizonte de 60 dias.}
	\begin{tabular}{|M{2.6cm}|c|c|c|c|c|c|c|c|c|c|}
\hline
\textbf{Inst\^ancia} & \textbf{M\'ed.} & \textbf{Unid.} & \textbf{$\Sigma$ Dem.} & \textbf{Dem./dia} & \textbf{$\Sigma$ Disp.} & \textbf{Disp./dia} & \textbf{$\Sigma$ Salas$\times$T} & \textbf{Salas/dia} & \textbf{Dem./Oferta} & \textbf{Disp./Dem.} \\
\hline\hline
02/01/22 a 02/03/22 & 146 & 35 & 129.617 & 2160 & 136.118 & 2269 & 5.784 & 96{,}4 & 22{,}41 & 1{,}05 \\
\hline
01/02/22 a 01/04/22 & 149 & 36 & 147.997 & 2467 & 151.902 & 2532 & 6.333 & 105{,}5 & 23{,}37 & 1{,}03 \\
\hline
01/04/22 a 30/05/22 & 162 & 36 & 147.217 & 2454 & 149.008 & 2483 & 6.421 & 107{,}0 & 22{,}93 & 1{,}01 \\
\hline
01/06/22 a 30/07/22 & 169 & 36 & 141.915 & 2365 & 149.923 & 2499 & 6.571 & 109{,}5 & 21{,}60 & 1{,}06 \\
\hline
01/08/22 a 29/09/22 & 160 & 37 & 152.426 & 2540 & 158.342 & 2639 & 6.843 & 114{,}0 & 22{,}27 & 1{,}04 \\
\hline
01/10/22 a 29/11/22 & 158 & 37 & 148.527 & 2475 & 152.139 & 2536 & 6.542 & 109{,}0 & 22{,}70 & 1{,}02 \\
\hline
01/12/22 a 29/01/23 & 175 & 37 & 133.009 & 2217 & 136.052 & 2268 & 5.931 & 98{,}8 & 22{,}43 & 1{,}02 \\
\hline
01/02/23 a 01/04/23 & 184 & 37 & 174.722 & 2912 & 171.118 & 2852 & 7.552 & 125{,}9 & 23{,}14 & 0{,}98 \\
\hline
01/04/23 a 30/05/23 & 193 & 37 & 177.365 & 2956 & 170.870 & 2848 & 7.489 & 124{,}8 & 23{,}68 & 0{,}96 \\
\hline
01/06/23 a 30/07/23 & 193 & 37 & 179.589 & 2993 & 177.056 & 2951 & 7.905 & 131{,}8 & 22{,}72 & 0{,}99 \\
\hline
\end{tabular}

	\label{quad:apendice-base-h60}
	\fautor
\end{sidewaysquadro}
\endgroup

\clearpage

\section{Horizonte de 90 dias}
\label{apendice:base-h90}

O Quadro~\ref{quad:apendice-base-h90}, ao final desta seção, apresenta a caracterização estendida das 10 instâncias trimestrais. Em relação ao Quadro~\ref{quad:base-h90} do capítulo, são acrescentadas as colunas de médias diárias (demanda, disponibilidade e salas\,$\times$\,turnos) e o número de unidades clínicas ativas no recorte.

\noindent No horizonte trimestral, cada instância corresponde a uma janela operacional completa de planejamento primário, com soma de exames demandados entre $\sim$205k e $\sim$273k. A média diária de demanda evolui de $\sim$2.284 (primeiro trimestre de 2022) para $\sim$3.030 (recortes de meados de 2023), acompanhada pela ampliação do quadro médico ativo (de 154 para 198 profissionais). As razões \emph{oferta\,$\times$\,demanda} apresentam a menor amplitude entre os quatro horizontes (\emph{Demanda / (Salas\,$\times$\,Turnos)} entre 21{,}9 e 23{,}6, \emph{Disponibilidade / Demanda} entre 0{,}97 e 1{,}05), refletindo o efeito de média do horizonte mais longo. Os $\sim$10\textsuperscript{6} variáveis binárias geradas por uma instância deste horizonte (cf.~Seção~\ref{cap:resultados:exato-gap}) impedem a convergência do método exato no limite de tempo adotado, motivando a comparação com as MIP-heurísticas e meta-heurísticas conduzida no Capítulo~\ref{cap:resultados}.

\begingroup\let\cleardoublepage\clearpage
\begin{sidewaysquadro}[p]
	\centering
	\caption{Caracterização estendida das 10 instâncias do horizonte de 90 dias.}
	\begin{tabular}{|M{2.6cm}|c|c|c|c|c|c|c|c|c|c|}
\hline
\textbf{Inst\^ancia} & \textbf{M\'ed.} & \textbf{Unid.} & \textbf{$\Sigma$ Dem.} & \textbf{Dem./dia} & \textbf{$\Sigma$ Disp.} & \textbf{Disp./dia} & \textbf{$\Sigma$ Salas$\times$T} & \textbf{Salas/dia} & \textbf{Dem./Oferta} & \textbf{Disp./Dem.} \\
\hline\hline
02/01/22 a 01/04/22 & 154 & 36 & 205.550 & 2284 & 213.625 & 2374 & 9.009 & 100{,}1 & 22{,}82 & 1{,}04 \\
\hline
01/02/22 a 01/05/22 & 162 & 36 & 218.818 & 2431 & 222.837 & 2476 & 9.439 & 104{,}9 & 23{,}18 & 1{,}02 \\
\hline
01/04/22 a 29/06/22 & 175 & 36 & 218.683 & 2430 & 224.851 & 2498 & 9.728 & 108{,}1 & 22{,}48 & 1{,}03 \\
\hline
01/06/22 a 29/08/22 & 174 & 37 & 217.185 & 2413 & 227.748 & 2531 & 9.910 & 110{,}1 & 21{,}92 & 1{,}05 \\
\hline
01/08/22 a 29/10/22 & 162 & 37 & 228.215 & 2536 & 236.405 & 2627 & 10.217 & 113{,}5 & 22{,}34 & 1{,}04 \\
\hline
01/09/22 a 29/11/22 & 162 & 37 & 222.336 & 2470 & 228.993 & 2544 & 9.909 & 110{,}1 & 22{,}44 & 1{,}03 \\
\hline
01/11/22 a 29/01/23 & 180 & 37 & 208.015 & 2311 & 212.281 & 2359 & 9.197 & 102{,}2 & 22{,}62 & 1{,}02 \\
\hline
01/01/23 a 31/03/23 & 189 & 37 & 246.456 & 2738 & 244.418 & 2716 & 10.724 & 119{,}2 & 22{,}98 & 0{,}99 \\
\hline
01/03/23 a 29/05/23 & 198 & 37 & 267.207 & 2969 & 257.980 & 2866 & 11.323 & 125{,}8 & 23{,}60 & 0{,}97 \\
\hline
01/05/23 a 29/07/23 & 195 & 37 & 272.689 & 3030 & 267.385 & 2971 & 11.851 & 131{,}7 & 23{,}01 & 0{,}98 \\
\hline
\end{tabular}

	\label{quad:apendice-base-h90}
	\fautor
\end{sidewaysquadro}
\endgroup

\clearpage

\section{Síntese}
\label{apendice:base-sintese}

A leitura conjunta dos quatro horizontes evidencia três padrões consistentes ao longo do conjunto histórico. Primeiro, o crescimento gradual do volume operacional --- demanda, disponibilidade e quantidade de médicos --- progride de forma aproximadamente linear entre 02/01/2022 e 31/07/2023, sem rupturas estruturais identificáveis nos quadros. Segundo, a razão \emph{Disponibilidade / Demanda} migra de um regime de leve excesso ($>1$) em 2022 para um regime de leve déficit ($<1$) em 2023, indicando que a expansão da força de trabalho registrada na base não acompanhou integralmente a ampliação da demanda no último semestre coberto. Terceiro, a amplitude das duas razões \emph{oferta\,$\times$\,demanda} diminui monotonicamente com o aumento do horizonte (de [20{,}8;\,23{,}5] em 15 dias para [21{,}9;\,23{,}6] em 90 dias na razão \emph{Demanda / (Salas\,$\times$\,Turnos)}), confirmando o efeito esperado de amortecimento estatístico em janelas mais longas. \textbf{Reitera-se} que estas razões mantêm o caráter de \emph{limitantes ingênuos} discutido na Seção~\ref{cap:resultados:base-geral}: oferecem um diagnóstico de regime, não uma medida de factibilidade.
